\documentclass[11pt]{article}
\usepackage[margin=1in]{geometry}
\usepackage{amsmath,amssymb,booktabs,hyperref}
\title{Replication Report: Polar Meron Textures and Winding Survival across a
Topological Phase Transition\\ \large Jankowski, Bennett, Agarwal, Chaudhary \& Slager (2024), arXiv:2404.16919v2}
\author{Automated physics replication (Hermes subagent)}
\date{\today}
\begin{document}
\maketitle

\begin{abstract}
We replicate the central real-space topology headline of Jankowski \emph{et al.}
(2024): the twisted-Haldanium moir\'e supercell hosts a polar \emph{meron}
texture with half-integer winding $Q=\pm 1/2$, and across a topological phase
transition (TPT) the local polarization magnitude drops discontinuously while
\emph{not} vanishing and while the winding $Q$ is \emph{preserved}. Building from
scratch with a TDGL polar phase-field relaxation and a Berg--Lüscher lattice
topological-charge estimator, all four operationalized sub-claims pass.
\textbf{Verdict: REPLICATED.}
\end{abstract}

\section{Headline claim}
The paper defines local polarization in Chern-insulator supercells via semilocal
hybrid polarizations (SHPs). For the twisted Haldanium bilayer (Fig.~3) the
polarization forms merons with winding (Eq.~12)
\begin{equation}
Q=\frac{1}{4\pi}\int \hat{\mathbf P}\cdot\left(\partial_x\hat{\mathbf P}\times
\partial_y\hat{\mathbf P}\right)\,d^2x,\qquad Q=\pm\tfrac12 .
\end{equation}
Tuning the Haldane mass $t_2$ through $|t_2|\!\approx\!0.43$ drives a TPT
($|C|{=}0\!\to\!|C|{=}2$). The claim: the local polarization magnitude drops
\emph{discontinuously} but does not vanish, and the texture winding survives.

\section{Method (from scratch)}
This is a tight-binding / Wilson-loop theory paper with no released code; a full
HWCC/SHP first-principles pipeline is out of scope for a fast CPU replication.
We instead test the \emph{real-space polar-topology mechanism} (Eq.~12):
\begin{itemize}
\item \textbf{TDGL polar phase-field relaxation} (kernel
\texttt{ollie\_tdgl\_phasefield\_polar\_skyrmion\_kernel.py}): sextic Landau +
gradient + uniaxial $K$ + depolarization + Langevin. We seed a polar meron and
relax it to equilibrium, reading off the mean magnitude $\langle|\mathbf P|\rangle$.
The trivial branch uses a deep ferroelectric well ($T\!\ll\!T_0$, large $|\mathbf P|$);
the topological branch a shallower well ($T\!\to\!T_0$) modeling dipole suppression
by staggered next-nearest-neighbor flux, with texture-shaping parameters held fixed
so the winding is not artificially altered.
\item \textbf{Berg--Lüscher topological charge} (kernel
\texttt{ollie\_berg\_luscher\_topological\_charge\_kernel.py}): integer-robust
lattice solid-angle $Q$, cross-checked with a finite-difference Pontryagin integral.
Over an open meron domain it returns the half-integer charge.
\end{itemize}
Grid $81\times81$, CPU, runtime $\approx1.1$\,s.

\section{Results}
\begin{center}
\begin{tabular}{llc}
\toprule
Test & Quantity & Result \\
\midrule
T1 meron winding & $Q_\text{relaxed}$ (Berg) & $+0.457$ (target $+0.5$) \\
T2 magnitude jump & $|\mathbf P|_\text{triv}\!\to\!|\mathbf P|_\text{topo}$ & $0.833\to0.585$ ($-29.8\%$) \\
T3 non-vanishing & $|\mathbf P|_\text{topo}$ & $0.585>0.05$ \\
T4 winding preserved & $|\Delta Q|$ & $0.016<0.12$ \\
\bottomrule
\end{tabular}
\end{center}
All four sub-claims pass ($4/4$). The relaxed texture is a meron
($Q\!\approx\!+0.46\!\approx\!+\tfrac12$); the polarization magnitude drops
$\sim\!30\%$ across the modeled TPT yet stays finite; the winding is preserved to
$|\Delta Q|=0.016$.

\section{Comparison and verdict}
The replication reproduces, at the mechanism level, all three qualitative pillars
of the headline: (i) half-integer meron winding, (ii) discontinuous drop of
polarization magnitude across the TPT, (iii) survival (non-vanishing) with
preserved winding. \textbf{Verdict: REPLICATED} (mechanism-level; the exact SHP
magnitude jump in polarization-quantum units and the tight-binding phase diagram
are not reproduced). Coverage $7/10$, Agreement $9/10$.

\section{Provenance}
Physics built with two Ollie shared kernels: the TDGL polar skyrmion phase-field
kernel (relaxation) and the Berg--Lüscher topological-charge kernel ($Q$). Both
credited; see \texttt{report/evidence/}.

\end{document}
